\section{After publication: the build of 2026-11-11 (cg31)}\label{sec:cg31}

Every table before this section is of snapshot cg29 or earlier. On 2026-11-11, the day after the repository was
published, twelve commits followed (patches 0028 to 0039): one on the energy schedule of a run without an end criteria
(0028), ten on the Metal diamond engine (0029 to 0038) and one on the cylinder extension (0039). The build, cg31, has
been the installed one on the author's Mac since that day. This section records what was looked for, what changed and
what it gave (\code{results/metal\_engine\_cg31\_2026-11-11.txt}).

\subsection{Where the Metal engine's time went}
Three observations set the work list.

\emph{An energy sum on one thread.} The fork's base estimates the field energy every Nyquist period during the first
60\,s of a run, so that a run with an end criteria stops at a reproducible step. The estimate is one pass of one CPU
thread over both fields, and the engine waits for it. On a benchmark model with a Nyquist period of 51 steps this was
4--6\,\% of the time-stepping loop on 3--17\,M cells; with a period of 519 steps the time lost outside the engine fell
from 0.23\,s to 0.03\,s per 3\,000 steps (project log). A run with a fixed number of steps and no end criteria has
nothing to stop on, so it now makes no such estimates (patch 0028), as in v0.37.0. The author's board runs have a
criteria and Nyquist periods above a thousand steps: for them the estimate costs 0.1--0.2\,\% (project log) and stays.

\emph{Extensions cost by being present.} In a timing-only build that leaves parts of the kernel out (results wrong by
design), the GCPW line ran 1.35$\times$ as fast without its Mur faces, 1.32$\times$ without its conducting sheets and
2.0$\times$ without any extension; the SMA launch 1.69$\times$ without its UPML slab on one z face.
But on a later snapshot, where the same builds of the GCPW line gave 1.26$\times$, 1.22$\times$ and 1.94$\times$,
skipping one stage at run time, with the code still in the kernel, gave only 1.02--1.07$\times$ each (the sheets' Pre
and Apply, Mur's Pre and Post). The cost is tied to the code being compiled into the kernel, not
to the loops over the few cells it handles; what in the GPU's execution makes it so was not found. Empty extra barriers
cost nothing measurable, and moving the Pre stages of Mur and of the sheets into the half-step before, beside the cell
update, which saves one barrier per half-step with identical fields, ran 2--4\,\% \emph{slower} and was dropped.

\emph{Not at the bus on mid-size grids.} Counted from the tile schedule, a tile that stays cached for its block loads
and stores 20--29 bytes of field per cell and time step, which at the measured rates is 79--105\,GB/s at 3\,M cells,
below the 125--138\,GB/s that STREAM measures on this machine (Section~\ref{sec:common}); the count assumes perfect
reuse within a block and was not measured. At 17\,M cells with 257 z lines a tile is 0.9--1.5\,MB, because a tile
holds every z line, and the shallower $K=3$ tilings win although they move more bytes (the byte counts and tile sizes
of this paragraph are from the project log).

\subsection{What changed}
\begin{table}[H]\centering\small
\begin{tabular}{@{}lp{0.60\textwidth}l@{}}
\toprule
patch & change & switch (default on) \\
\midrule
0028 & no energy schedule in a run without an end criteria & \\
0029, 0032 & Mur faces of one axis share a barrier; Post and Apply in one pass where a tile's faces are on one axis &
  \code{DIAMOND\_MURMERGE} \\
0029 & a call's rest in blocks of nearly equal depth (10 steps at $K=4$: 4, 3, 3) & \code{DIAMOND\_EVEN} \\
0030 & three dispatches per block: the mountain$\times$valley and valley$\times$mountain tiles do not depend on each
  other and go together & \code{DIAMOND\_MERGE} \\
0031, 0032, 0034 & no divergent branch for the z neighbour; divisions by run-time widths as multiplications by
  reciprocals from the host; threads without work leave a Mur stage early & \code{DIAMOND\_FASTIDX} \\
0031 & the pipeline is told that a threadgroup is a whole number of SIMD groups & \code{DIAMOND\_TGFLAG} \\
0033 & UPML of the boxes with every z line (the x and y slabs) a z-vector at a time & \code{DIAMOND\_PML4} \\
0035 & UPML of the z slabs three components of a cell at a time & \code{DIAMOND\_PML3} \\
0036 & dispersive materials (sheets, Drude/Lorentz, Debye) three components of a cell at a time & \\
0037 & offsets and counts of the extension data as constants of the kernel, which is compiled after the data are
  packed & \\
0038 & seven tilings in a two-stage tuner (adds $K=4$ with 10,6 and 9,2 tiles and $K=5$ with 12,8); the per-tile
  lists in their own buffer & \code{DIAMOND\_CAND} \\
0039 & line hooks for the cylinder extension (Section~\ref{sec:cg31cyl}) & \\
\bottomrule
\end{tabular}
\caption{The commits of 2026-11-11. Switch names carry the prefix \code{OPENEMS\_}; \code{=0} restores the earlier
behaviour (\code{DIAMOND\_CAND=4}: the four tilings of patch 0016). Every step was a snapshot of its own (cg30a to
cg30l, cg31a) that passed \code{scripts/checks/diamond\_check.sh} with 72 of 72 comparisons
identical.}\label{tab:cg31items}
\end{table}

The mountain$\times$valley and the valley$\times$mountain tiles of a block can share a dispatch because a mountain
cell reads only cells that were mountain cells one half-step earlier, in both directions: a cell leaves its x mountain
or its y mountain first, so it is updated by one of the two kinds only, and neither kind reads what the other writes
within the block. The UPML stages were scalar loops over single floats, three components and up to four lanes apart;
an x or y slab has every z line, so its cells fill whole z-vectors and one pass can update a float4 of one component
with the same operations per lane, while a z slab has one lane of a z-vector per cell and its pass now takes the
three components of a cell together. One build in this series was wrong: it computed the reciprocals for the
divisions on the GPU, and on this M5 (macOS 27.0.1) the GPU's own \code{0xffffffff / w} returns 3\,579\,139\,454 for
$w=3$, where 1\,431\,655\,765 is right. The identity check caught it as wrong fields on the Mur cases; the
reciprocals now come from the host.

\subsection{What it gave}
\begin{table}[H]\centering\small
\begin{tabular}{@{}lrrrrrr@{}}
\toprule
 & \multicolumn{3}{c}{Metal diamond engine} & \multicolumn{2}{c}{CPU engine, cg31} & Metal / \\
\cmidrule(lr){2-4}\cmidrule(lr){5-6}
model & cg29 & cg31 & cg31 / cg29 & 8 threads & 4 threads & CPU best \\
\midrule
GCPW line, 0.24\,M & 2\,118 & 2\,307 & 1.09$\times$ & 2\,000 & 1\,816 & 1.15$\times$ \\
SMA launch, 0.87\,M & 1\,854 & 2\,108 & 1.14$\times$ & 1\,955 & 1\,509 & 1.08$\times$ \\
patch tutorial, 0.10\,M & 1\,753 & 1\,887 & 1.08$\times$ & 1\,910 & 1\,868 & 0.99$\times$ \\
notch tutorial, 0.27\,M & 1\,904 & 2\,907 & 1.53$\times$ & 2\,129 & 1\,909 & 1.37$\times$ \\
patch x5, 7.86\,M & 2\,901 & 3\,100 & 1.07$\times$ & 2\,104 & 1\,802 & 1.47$\times$ \\
notch x3, 5.25\,M & 3\,210 & 3\,702 & 1.15$\times$ & 2\,330 & 1\,853 & 1.59$\times$ \\
\addlinespace
PEC box, 3.0\,M & 3\,692 & 4\,408 & 1.19$\times$ & 2\,614 & 2\,156 & 1.69$\times$ \\
PEC box, 17\,M & 3\,056 & 3\,507 & 1.15$\times$ & 1\,736 & 1\,550 & 2.02$\times$ \\
cube with 10-cell PML, 1.0\,M & 980 & 1\,546 & 1.58$\times$ & 1\,343 & 1\,012 & 1.15$\times$ \\
cube with 10-cell PML, 8.0\,M & 1\,483 & 1\,839 & 1.24$\times$ & 1\,293 & 1\,049 & 1.42$\times$ \\
\bottomrule
\end{tabular}
\caption{The Metal diamond engine before (cg29) and after (cg31) the commits of 2026-11-11, with cg31's CPU engine
beside it: median \MC\ of four paired rounds per model on 2026-11-11 with the other sessions on the machine paused
(the other processes used 0.2--0.9 cores; a run that saw more than 1.5 cores was repeated once). cg31 / cg29: the
median of the per-round ratios, which over the rounds stayed within 3\,\% of the value shown except on the patch
tutorial (0.92--1.13), the notch tutorial (1.42--1.53) and the two cubes (1.51--1.58 and 1.20--1.29). The PEC box is
the benchmark model of luikore's engine (Table~\ref{tab:threebuilds}), the cube that of Table~\ref{tab:cube}. The SMA
launch is on the toolkit's current mesh. Every cg31 Metal run's probes are identical to cg29's.}\label{tab:cg31}
\end{table}

The two boards gain 9 and 14\,\%; the notch tutorial and the PML cube gain the most (53\,\%, and 58 and 24\,\%), in
good part from the UPML stages on whole faces (patch 0033); the two large tutorial models gain 7 and 15\,\% and the
PEC box 19 and 15\,\%. With that the Metal engine is ahead of the CPU
engine on every model but the patch tutorial, where the two are level. cg31 also passed, on the CPU side,
\code{emskit fdtd check} on the six models against stock (all identical, 18 of 18 rows), the Debye check and both
boards end to end, which stop at the same steps as with cg29 (34\,300 and 50\,320); the validation suite on the Metal
engine (15 of 15) ran on the snapshot of patches 0001--0038 (cg30l), before the cylinder patch.

\subsection{Three builds on one model}\label{sec:threebuilds}
Section~\ref{sec:gpu} had no head-to-head with luikore's Metal engine~\cite{luikore228}: its diamond path refuses the
six models. It does run its own benchmark model, a uniform grid with PEC walls, one dielectric box and a point
excitation, so the three builds were run on that model, each binary from the same model file, interleaved, with the
other sessions paused (Table~\ref{tab:threebuilds}; \code{scripts/luikore\_compare/},
\code{results/engines\_on\_luikore\_model\_2026-11-11.txt}).

\begin{table}[H]\centering\small
\begin{tabular}{@{}lrr@{}}
\toprule
engine & 3.0\,M cells & 17.0\,M cells \\
\midrule
fdtdKit (cg31), Metal diamond & 3\,981 & 3\,309 \\
luikore PR \#228 (55ce11b), Metal & 3\,952 & 2\,829 \\
master 1d710c2, \code{gpu} & 1\,638 & 1\,566 \\
master 1d710c2, \code{cpu}, 8 threads & 1\,872 & 1\,933 \\
fdtdKit (cg31), \code{multithreaded}, best of 8 and 10 threads & 2\,371 & 1\,659 \\
luikore PR \#228, \code{multithreaded}, 8 threads & 1\,296 & 1\,042 \\
master 1d710c2, \code{multithreaded}, 8 threads & 1\,246 & 1\,055 \\
\bottomrule
\end{tabular}
\caption{Three builds on the benchmark model of luikore's engine, one Apple M5, 2026-11-11: median \MC\ of three
interleaved repetitions (3\,000 steps at 3.0\,M cells, 1\,000 at 17.0\,M); the other processes used a median of 0.44
cores before a run, at most 1.14. The \code{multithreaded} engines of the three builds give identical probes; the
others differ from them by $7\times10^{-6}$ to $3\times10^{-5}$ relative.}\label{tab:threebuilds}
\end{table}

On its own model luikore's engine and cg31's are level at 3.0\,M cells (1.01$\times$) and cg31's is 1.17$\times$ as
fast at 17\,M; the night before, cg29's had run 3\,728 and 3\,092, 0.94$\times$ and 1.09$\times$ luikore's in that
run (3\,958 and 2\,848; \code{results/luikore\_compare\_cg29\_2026-11-11.txt}). cg31's engine is 2.4$\times$ and
2.1$\times$ master's \code{gpu} engine here, luikore's 2.4$\times$ and 1.8$\times$.
On the CPU, cg31 is 1.27$\times$ master's new \code{cpu} engine at 3.0\,M cells and 0.86$\times$ at 17\,M: on the large
grid master's engine is ahead, as in Table~\ref{tab:upstream}. Two things qualify the table. It is one model, the best
case of a GPU engine (a regular grid, no boundary work). And a model file whose \code{endCriteria} is 0 makes the
\code{openEMS} program use its default of $-60$\,dB, so here the fork and master make the energy estimates of
Section~\ref{sec:cg31} and luikore's build does not: run from Python, where 0 means no criteria, cg31's Metal engine
does 4\,408 and 3\,507\,\MC\ on this model (Table~\ref{tab:cg31}).

Master's \code{cpu} engine blocks in time: all threads work in one tile of the size of the performance cores' cache
for up to eight steps. With that switched off (\code{OPENEMS\_CPU\_TBLOCK=0}) it ran upstream's ideal case at 8\,M
cells at 1\,347 and 1\,436\,\MC\ (4 and 8 threads) against 2\,010 and 2\,099 blocked the day before, and the PML case
at 891 and 1\,069 against 1\,117 and 1\,225: the blocking is worth 1.46--1.49$\times$ and 1.15--1.25$\times$ on this
machine (two runs a day apart). The fork's fused sweep, at 1\,739 on the ideal case, is ahead of master's unblocked
kernel and behind its blocked one. The experiment of Section~\ref{sec:rejected} that found no gain in temporal blocking
on the CPU gave every thread a sub-grid of its own, which does not fit the cache four threads share; one tile for
all threads is the next CPU item.

\subsection{A closed cylinder on the fused sweep}\label{sec:cg31cyl}
The cylinder extension, which closes a mesh in $\alpha$ and handles $r=0$, declared no hooks, so such a mesh ran the
unfused plan with the extension's two Post hooks on one thread. Both hooks stay on the x-line (r-line) they read and
keep no state, so they are line hooks now (patch 0039) and the mesh runs the fused sweep with the neighbour flags. On
a closed cylinder of 61$\times$72$\times$41 lines with $r=0$ the engine ran 1\,857 and 1\,891\,\MC\ at 8 threads before
and 2\,503 and 2\,553 after (two rounds in the quiet block; at 4 threads the rounds disagree, 2\,126 and 2\,176
against 2\,125 and 2\,400; stock 464--474). Five probes on the axis and the $\alpha$ seam equal stock's, the raw
fields are bit-identical to the unfused plan and at 1, 4, 8 and 10 threads, with and without $r=0$, and the
cylindrical multigrid tutorial, which stays unfused, has 18 datasets with no difference
(\code{scripts/cylinder\_closed/}). The first version of this test had its source between two mesh lines: every
probe was zero and every comparison passed. The script now puts the source and the probes on mesh lines.

\subsection{A shorter pulse (tested, not adopted)}\label{sec:pulse}
The toolkit excites with a Gaussian of centre $f_{\max}/2$ and width $f_{\max}/2$. A baseband Gaussian (centre 0,
width $f_{\max}$) covers the same band with half the length, and a toolkit run lasts the pulse plus a hold of half a
pulse length. Both boards were run with both pulses and compared with the toolkit pulse run to the full time
(Table~\ref{tab:pulse}; \code{scripts/pulse/}, \code{results/pulse\_test\_2026-11-11.txt}; the excitation was set on
the built model, the toolkit was not changed).

\begin{table}[H]\centering\small
\begin{tabular}{@{}llrrrr@{}}
\toprule
board & pulse, hold & steps & max $\Delta S_{21}$ & max $\Delta|S_{11}|$ & max $\Delta Z_0$ \\
\midrule
GCPW line & toolkit, 0.5 & 34\,300 & 0.0007\,dB & $1.0\times10^{-4}$ & 0.05\,$\Omega$ \\
 & baseband, 0.5 & 22\,540 & 0.0042\,dB & $5.5\times10^{-4}$ & 0.09\,$\Omega$ \\
 & baseband, 1.0 & 27\,930 & 0.0016\,dB & $2.2\times10^{-4}$ & 0.17\,$\Omega$ \\
SMA launch & toolkit, 0.5 & 50\,320 & 0.0003\,dB & $2.1\times10^{-5}$ & 0.66\,$\Omega$ \\
 & baseband, 0.5 & 25\,900 & 0.0004\,dB & $4.8\times10^{-5}$ & 0.53\,$\Omega$ \\
 & baseband, 1.0 & 34\,040 & 0.0005\,dB & $4.2\times10^{-5}$ & 0.51\,$\Omega$ \\
\bottomrule
\end{tabular}
\caption{The toolkit's pulse and a baseband pulse of half its length on the two boards, probe end at $-60$\,dB with
the hold in lengths of the pulse used, against the toolkit pulse run to the full time (94\,217 and 118\,475 steps),
50\,MHz to 2\,GHz. $|S_{11}|$ linear: the GCPW line's $S_{11}$ is $-40$ to $-52$\,dB, the SMA launch's $-17$ to
$-23$\,dB.}\label{tab:pulse}
\end{table}

The baseband pulse cuts the steps 1.52$\times$ on the GCPW line and 1.94$\times$ on the SMA launch; the errors grow a
few-fold and stay small ($S_{21}$ within 0.005\,dB, $|S_{11}|$ at the $-65$\,dB level and below). It is not the
toolkit's default: a baseband pulse carries DC, and a board with a conductor that has no DC path, or a series
capacitor, was not tested; there the signals may not settle and the run would go to its full time.
